\documentclass[12pt]{article}
\usepackage{setspace}
\usepackage{amsmath,amssymb,euscript,graphicx}
\usepackage{xurl}
\usepackage{graphicx}
\usepackage{subfigure}
\usepackage{wrapfig}
\usepackage[compact]{titlesec}
\usepackage{fancybox,color}
\usepackage[footnotesize]{caption}
\usepackage{cite}
\usepackage{enumitem}
\usepackage[normalem]{ulem}
\usepackage[top=0.9375in, right=0.9375in, left=0.9375in, bottom=0.9375in, centering]{geometry}
\usepackage{bm}
\usepackage{float}

\usepackage{amsthm}


% -- Loading the code block package:
\usepackage{listings}
% -- Basic formatting
\usepackage[utf8]{inputenc}
\usepackage[english]{babel}
\usepackage{times}
\setlength{\parindent}{8pt}
\usepackage{indentfirst}
% -- Defining colors:
\usepackage[dvipsnames]{xcolor}
\definecolor{codegreen}{rgb}{0,0.6,0}
\definecolor{codegray}{rgb}{0.5,0.5,0.5}
\definecolor{codepurple}{rgb}{0.58,0,0.82}
\definecolor{backcolour}{rgb}{0.95,0.95,0.92}
% Definig a custom style:
\lstdefinestyle{mystyle}{
    backgroundcolor=\color{backcolour},   
    commentstyle=\color{codepurple},
    keywordstyle=\color{NavyBlue},
    numberstyle=\tiny\color{codegray},
    stringstyle=\color{codepurple},
    basicstyle=\ttfamily\footnotesize\bfseries,
    breakatwhitespace=false,         
    breaklines=true,                 
    captionpos=t,                    
    keepspaces=true,                 
    numbers=left,                    
    numbersep=5pt,                  
    showspaces=false,                
    showstringspaces=false,
    showtabs=false,                  
    tabsize=2
}
% -- Setting up the custom style:
\lstset{style=mystyle}


% IEEE uses Times Roman font, so we'll default to Times.
% These three commands make up the entire times.sty package.
\renewcommand{\sfdefault}{phv}
\renewcommand{\rmdefault}{ptm}
\renewcommand{\ttdefault}{pcr}

\newcommand{\squeeze}{\vspace{-0.1in}}
\newcommand{\changed}[1]{\textcolor{blue}{#1}} %generates text in blue

%%%%%%%%%%%%%%%%%%%%%%%%%%%%%%%%%%%%%%%%%%%%%%%%%%%%%%%%%%%%%%%%%%%%%%%

\begin{document}

\begin{center}
  {\Large {\bf Lab 2: Functions that Control Shapes}}
  
\end{center}

\vspace{5pt}
\hrule
\vspace{5pt}

\section{Summary}
In this lab, we will add further generality to our shapes using functions and parameters.
The main goal is to learn about good code organization. 

\section{Setup}
\begin{itemize}
  \item open a terminal in VS Code. 
  \item Navigate to your ES1098 folder, use \texttt{ls}, \texttt{cd}, and \texttt{pwd} to help you find this directory
  \item Once you are in this directory you should see: \texttt{/Users/<computername>/ES1098} when you type \texttt{pwd}.
  \item Create a new folder by typing \texttt{mkdir lab2}
  \item Move into that folder by typing \texttt{cd lab2}
  \item Confirm your current directory by typing \texttt{pwd}.
\end{itemize}

In VS Code, make a new file called \texttt{lab2.py}. Remember: you can do this using the \texttt{code lab2.py} in the terminal of your VS Code IDE. Make sure you are in the correct directory. 

\section{Code Organization}
In \texttt{lab2.py} write a function that draws a triangle. It should take in one parameter, \texttt{sideLength}, which is the length (in pixels) of each of the sides. \textbf{Important}: Draw the triangle so that the turtle starts out at the lower-left corner. Remember the turtle initiall faces eastward on the screen.

Call the triangle function from main code, or the code not indented at the bottom of your file. Test your code by running it in the terminal \textit{python lab2.py}.

Now we are going toconsider the code organization and good practices.

\begin{enumerate}
\item Header comments: which includes your name, course, semester, and a one line description
\item Import statements, \texttt{import turtle}
\item Function definitions separated by empty lines
\item Every funciton should have a comment that describes what the functio ndoes
\item At least one inline comment if there are more than 3 lines of code in the function
\item A comment indeication where the main code starts
\item Main code the program will run
\end{enumerate}
Check out the turtle.speed() function to see how you can speed up the drawing of your figures.
Don't forget to use the \texttt{exitonclick()} functionat the end of your file 


\section{Moving the turtle anywhere on the screen}
The goal of this section is to draw a triangle anywhere on the canvas, not just at the (0, 0) position. 

\section{Positioning and scaling shapes}
We want to draw shapes at any size (also called scale) and position.
\subsection{Draw a triangle at any position and scale}
Write a new function called \texttt{triangle2} hat has the following signature:

\begin{lstlisting}[language=Python]
def triangle2(x, y, scale):
    """
    Draws a triangle at any (x, y) position and scale (scale).
    By default, the side Lengths are 100 (when scale = 1)
    """
    # Your code goes here
\end{lstlisting}
Your \texttt{triangle2} function should move the turlte to the given (x, y) position and draw a triangle with a side length equal to 100 pixels when the scale parameter is 1.

Guidelines:
\begin{itemize}
\item You should only need to write 2 lines of code
\item Remember you can call functions inside of other functions
\end{itemize}

\subsection{Testing your \texttt{triangle2} function}

Comment out any existing function calls in the main code that are not \texttt{speed}, or \texttt{exitonclick}. Test \texttt{triangle2} by callign it 3 times: with scale = 1 (standard size), scale = 2 (double standard size), and scale = 3 (triple standard size).
The output does not need to be fancy, the key is that your triangles appear at different positions and sizes. Here is an example: 

\begin{figure}[H]
\centering
\includegraphics[width=0.5\textwidth]{lab2_4n.png}
\end{figure}

\section{Position and scale compound shape}
We are going to make a compound shape, or a more complex shape made up of your simpler shapes arranged in different positions and / or scales.

\subsection{Draw a hard-coded compound shape}
Copy the following code into your \texttt{lab2.py} file. 
\begin{lstlisting}[language=Python]
def triangleStack(x, y, scale):
    '''Draws three triangles on top of each other.
    Smaller triangles are placed on top of larger ones
    '''
    # Largest triangle
    triangle2(0, 0, 2)
    # Medium triangle in middle
    triangle2(50, 173.2, 1)
    # Small triangle on top of stack
    triangle2(75, 259.81, 0.5)
\end{lstlisting}

Next, comment out existing main code except for \texttt{speed} and \texttt{exitonclick}.
Finally, call \texttt{triangleStack} to see what it looks like from the main code.
Note, the parameters are not used yet, so pick any 3 values right now when calling the function, \textit{e.g.}, \texttt{triangleStack(0, 0, 1)}

\subsection{Convert hard-coded positions to offsets}
We are going to convert the hard-coded compound shape (stack of three triangles) into a shape that we can move to any position and make any scale that we wish. We will do this in two stages.

First, convert the hard-coded positions to offsets. Make the lower-left corner of the first triangle (where the turtle starts out) "the origin" at (x, y), i.e., has offset of (0, 0). Compute the offsets relative to (x, y).

When you're done, test it by commenting out existing main code and draw \texttt{triangleStack} at three positions: \texttt{(x, y) = (-200, 0), (x, y) = (0, 0)}, and \texttt{(x, y) = (200, 0)}. The scale parameter value does not matter right now.

The output should look like this: 
\begin{figure}[H]
\centering
\includegraphics[width=0.5\textwidth]{lab2_5b.png}
\end{figure}

\subsection{Scaling the compound shape}
The second and last step is to make the triangle stack resizable. Do this by multiplying the offsets and scales by the parameter scale.

When you are done, test it by modifying your existing main code:
\begin{itemize}
\item Draw the 1st stack at scale 1, the 2nd at scale 1/2 and the third at scale 1/3
  \item Draw at these 3 positions texttt{(x, y) = (-200, 0), (x, y) = (0, 0)}, and \texttt{(x, y) = (100, 0)}
\end{itemize}

The ouput should look like this: 
\begin{figure}[H]
\centering
\includegraphics[width = 0.5\textwidth]{lab2_5c.png}
\end{figure}

Take a screen shot of your final product. This is \textbf{Required Image 1} for your project report.

\section{Acknowledgments}
This lab is from Professor Hannen Wolfe in the Colby College CS Department. 
Sources:
\begin{itemize}
\item \footnotesize{\url{https://cs.colby.edu/hewolfe/courses/cs151/f21/projects/02/Lab02FunctionsToControlShapes.html}}
\end{itemize}


\end{document}
